\resizebox{\textwidth}{!}{%
\begin{tikzpicture}[
    font=\small,
    >=Stealth,
    stage/.style={rectangle, rounded corners=2pt, draw=#1!70!black, fill=#1!10, minimum width=2.6cm, minimum height=0.95cm, align=center, line width=0.6pt},
    item/.style={rectangle, rounded corners=2pt, draw=#1!70!black, fill=#1!8, minimum width=2.15cm, minimum height=0.68cm, align=center, line width=0.5pt},
    data/.style={cylinder, shape border rotate=90, aspect=0.22, draw=#1!70!black, fill=#1!8, minimum width=1.95cm, minimum height=0.9cm, align=center, line width=0.5pt},
    arrow/.style={->, line width=0.75pt, draw=black!70},
    group/.style={rectangle, rounded corners=3pt, draw=black!35, fill=black!2, inner sep=8pt}
]

\node[item=teal] (pdf) {PDF/Word\\技术文档};
\node[item=teal, below=0.25cm of pdf] (img) {设备图像\\拓扑图};
\node[item=teal, below=0.25cm of img] (flow) {操作规程\\流程图};

\node[stage=cyan, right=1.05cm of img] (parse) {多模态解析\\与模态分流};
\node[stage=blue, right=1.15cm of parse] (extract) {LLM/VLM\\实体关系抽取};
\node[stage=violet, right=1.15cm of extract] (unify) {统一节点表示\\名称/类型/来源};
\node[stage=orange, right=1.15cm of unify] (embed) {BGE-M3\\语义嵌入};
\node[data=green, right=1.15cm of embed] (neo4j) {Neo4j\\知识图谱};
\node[data=red, below=0.6cm of neo4j] (index) {HNSW\\向量索引};

\draw[arrow] (pdf.east) -- (parse.west);
\draw[arrow] (img.east) -- (parse.west);
\draw[arrow] (flow.east) -- (parse.west);
\draw[arrow] (parse) -- (extract);
\draw[arrow] (extract) -- (unify);
\draw[arrow] (unify) -- (embed);
\draw[arrow] (embed) -- (neo4j);
\draw[arrow] (embed.south east) to[out=-25,in=180] (index.west);
\draw[arrow] (index.north) -- (neo4j.south);

\node[below=0.18cm of extract, align=center, text=black!65] (triples) {实体、关系、流程步骤\\统一转化为三元组};
\draw[arrow, dashed, draw=black!45] (extract.south) -- (triples.north);
\draw[arrow, dashed, draw=black!45] (triples.east) -- (unify.south);

\begin{scope}[on background layer]
    \node[group, fit=(pdf) (img) (flow), label={[font=\small]above:输入层}] {};
    \node[group, fit=(parse) (extract) (unify) (embed), label={[font=\small]above:构建层}] {};
    \node[group, fit=(neo4j) (index), label={[font=\small]above:存储与检索层}] {};
\end{scope}
\end{tikzpicture}%
}
